% problem_id: S001-lemma-approximate-solution
% source_id: S001 (Stacks Project)
% source_file: more-algebra.tex
% statement_locator: more-algebra.tex lines 35774--35788
% proof_locator: more-algebra.tex lines 35790--35890
% env: lemma
% pairing: tight (verified)
% status: complete (statement + author proof verbatim + reason + locators)
%
\begin{lemma}
\label{lemma-approximate-solution}
Let $p$ be a prime number. Let $A \subset B$ be an extension of discrete
valuation rings with fraction field extension $L/K$. Let $K_2/K_1/K$ be
a tower of finite field extensions. Assume
\begin{enumerate}
\item $K$ has characteristic $p$,
\item $L/K$ is separable,
\item $B$ is Nagata,
\item $K_2$ is a solution for $A \subset B$,
\item $K_2/K_1$ is purely inseparable of degree $p$.
\end{enumerate}
Then there exists a separable extension $K_3/K_1$
which is a solution for $A \subset B$.
\end{lemma}

\begin{proof}
Let us use notation as in Remark \ref{remark-construction}; we will use all
the observations made there. Since $L/K$ is separable, the algebra
$L_1 = L \otimes_K K_1$ is reduced
(Algebra, Lemma \ref{algebra-lemma-separable-extension-preserves-reducedness}).
Since $B$ is Nagata, the ring extension $B \subset B_1$ is finite
where $B_1$ is the integral closure of $B$ in $L_1$
and $B_1$ is a Nagata ring. Similarly, the ring $A$ is Nagata by
Lemma \ref{lemma-nagata-goes-down} hence $A \subset A_1$ is finite
and $A_1$ is a Nagata ring too. Moreover, the same assertions are true
for $K_2$, i.e., $L_2 = L \otimes_K K_2$ is reduced,
the ring extensions $A_1 \subset A_2$ and $B_1 \subset B_2$
are finite where $A_2$, resp.\ $B_2$ is the integral closure of
$A$, resp.\ $B$ in $K_2$, resp.\ $L_2$.

\medskip\noindent
Let $\pi \in A$ be a uniformizer. Observe that $\pi$ is a nonzerodivisor
on $K_1$, $K_2$, $A_1$, $A_2$, $L_1$, $L_2$, $B_1$, and $B_2$ and we have
$K_1 = (A_1)_\pi$, $K_2 = (A_2)_\pi$, $L_1 = (B_1)_\pi$, and $L_2 = (B_2)_\pi$.
We may write $K_2 = K_1(\alpha)$ where $\alpha^p = a_1 \in K_1$, see
Fields, Lemma \ref{fields-lemma-finite-purely-inseparable}.
After multiplying $\alpha$ by a power of $\pi$ we may and do assume
$a_1 \in A_1$. For the rest of the proof it is convenient to write
$K_2 = K_1[x]/(x^p - a_1)$ and $L_2 = L_1[x]/(x^p - a_1)$.
Consider the extensions of rings
$$
A'_2 = A_1[x]/(x^p - a_1) \subset A_2
\quad\text{and}\quad
B'_2 = B_1[x]/(x^p - a_1) \subset B_2
$$
We may apply Lemma \ref{lemma-construct-extension} to
$A'_2 \subset A_2$ and $f = \pi^2$ and to $B'_2 \subset B_2$ and $f = \pi^2$.
Choose an integer $n$ large enough which works for both of these.

\medskip\noindent
Consider the algebras
$$
K_3 = K_1[x]/(x^p - \pi^{2n} x - a_1)
\quad\text{and}\quad
L_3 = L_1[x]/(x^p - \pi^{2n} x - a_1)
$$
Observe that $K_3/K_1$ and $L_3/L_1$ are finite \'etale algebra extensions of
degree $p$. Consider the subrings
$$
A'_3 = A_1[x]/(x^p - \pi^n x - a_1)
\quad\text{and}\quad
B'_3 = B_1[x]/(x^p - \pi^n x - a_1)
$$
of $K_3 = (A'_2)_\pi$ and $L_3 = (B'_3)_\pi$. We are going to construct
a commutative diagram
$$
\xymatrix{
B'_2/\pi^{2n} B'_2 \ar[r]_{\psi'} & B'_3/\pi^{2n} B'_3 \\
A'_2/\pi^{2n} A'_2 \ar[r]^{\varphi'} \ar[u] & A'_3/\pi^{2n} A'_3 \ar[u]
}
$$
Namely, $\varphi'$ is the unique $A_1$-algebra isomorphism
sending the class of $x$ to the class of $x$.
Similarly, $\psi'$ is the unique $B_1$-algebra isomorphism sending
the class of $x$ to the class of $x$. By our choice of $n$ we obtain, via
Lemma \ref{lemma-construct-extension} and
Remark \ref{remark-functoriality-construct-extension}
finite ring extensions $A'_3 \subset A_3$ and $B'_3 \subset B_3$
such that $A'_3 \to B'_3$ extends to a ring map $A_3 \to B_3$
and a commutative diagram
$$
\xymatrix{
B_2/\pi^2 B_2 \ar[r]_\psi & B_3/\pi^2B_3 \\
A_2/\pi^2 A_2 \ar[r]^\varphi \ar[u] & A_3/\pi^2A_3 \ar[u]
}
$$
with all the properties asserted in the references mentioned above
(in particular $\varphi$ and $\psi$ are isomorphisms).

\medskip\noindent
With all of this data in hand, we can finish the proof. Namely, we
first observe that $A_3$ and $B_3$ are finite products of
Dedekind domains with $\pi$ contained in all of the maximal ideals.
Namely, if $\mathfrak p \subset A_3$ is a maximal ideal, then
$\pi \in \mathfrak p$ as $A \to A_3$ is finite. Then $\mathfrak p/\pi^2 A_3$
corresponds via $\varphi$ to a maximal ideal in $A_2 / \pi^2A_2$
which is principal as $A_2$ is a finite product of Dedekind domains.
We conclude that $\mathfrak p/\pi^2 A_3$ is principal and hence by Nakayama 
we see that $\mathfrak p (A_3)_\mathfrak p$ is principal.
The same argument works for $B_3$. We conclude that
$A_3$ is the integral closure of $A$ in $K_3$
and that $B_3$ is the integral closure of $B$ in $L_3$.
Let $\mathfrak q \subset B_3$ be a maximal ideal lying
over $\mathfrak p \subset A_3$. To finish the proof we have
to show that $(A_3)_\mathfrak p \to (B_3)_\mathfrak q$ is
formally smooth in the $\mathfrak q$-adic topology. By the criterion of
Lemma \ref{lemma-extension-dvrs-formally-smooth} it suffices to show that
$\mathfrak p (B_3)_\mathfrak q = \mathfrak q (B_3)_\mathfrak q$
and that the field extension $\kappa(\mathfrak q)/\kappa(\mathfrak p)$
is separable. This is true because we may check both assertions
by looking at the ring map $A_3/\pi^2 A_3 \to B_3/\pi^2 B_3$
and this is isomorphic to the ring map
$A_2/\pi^2 A_2 \to B_2/\pi^2 B_2$ where the corresponding
statement holds by our assumption that $K_2$
is a solution for $A \subset B$. Some details omitted.
\end{proof}
